\chapter{Аналитическая часть}
%<пишем про тему и все, связанное с ней>

Конвейерная обработка — это метод повышения производительности, при котором сложная задача разбивается на последовательность простых этапов, и разные этапы выполняются одновременно над разными частями данных или инструкций. 
Конвейер обработки данных состоит из трёх основных компонентов: источника, одного или нескольких этапов обработки и пункта назначения.~\cite{lit2}

\section{Задание}

Конвейер из k потоков, включая 3 отдельных на каждое ОУ и (k-3) дополнительно для ОУ2. k взять как рекомендованное значение для выбранной АЭВМ из ЛР4. 

Обслуживающее устройство 1 выполняет считывание графа в формате матрицы смежности из текстового или бинарного файла (имя файла дано в заявке). 

Обслуживающее устройство 2 выполняет обработку графа, согласно заданию на ЛР4, задействовав k-3 вспомогательных потоков. 

Обслуживающее устройство 3 выполняет дамп ответа в файл, имя которого получено добавлением префикса к имени входного файла из заявки (файлы с ответами соотносятся с заявками 1:1). 

3 нативных потока, которые передают друг другу информацию через очереди. Данные в 1-ю очередь (входную для 1-го обслуживающего устройства) в количестве N заявок заполнить перед пуском конвейера, считав их из файла (список имен файлов).

\section*{Вывод}

В данном разделе мной был проведен анализ теоретических материалов, объясняющих принцип работы алгоритмов.

\clearpage
